\lecture{4}{Mon 20 Feb 2023 16:33}{Properties of measure}

\begin{theorem}
\begin{itemize}
	\item 
	The Lebesgue measureble sets are closed in completmenet, countable union, finite intersection and set exclusion.
\item Countable additivity: if $E_n$ is pairwise disjoint measurable sets, then
	\[
		m\left( \bigcup_{n \in \mathbb{N}} E_n \right)  = \sum_{n=1}^{\infty} m(E_n)
	.\] 
\item every open set is measurable.
\end{itemize}
\end{theorem}

\begin{lemma}
	If $E, F \in \mathcal{M}$ and $E, F$ disjoint, then for any $A \subset \mathbb{R}$ we have
	\[
		m^*\left( A \cap (E \cup F) \right)  = m^* (A \cap E) + m^* (A \cap F) \qquad (*)
	.\] 
\end{lemma}
\begin{proof}
	\begin{itemize}
		\item Case 1: $A \subset E \cup F$. Then $A \setminus  E = A \cap F$, then the claim follows from definition of measurability:
			\[
				m^*(A) = m^*(A \cap E) + m^*(A \setminus E)
			.\] 
		\item For general $A$ : We replace $A $ by $A' = A \cap (E \cup F) \subset E \cup F$, and none of the terms in (*) changes.
	\end{itemize}
\end{proof}

\begin{proof}[Proof of theorem]
	\begin{itemize}
		\item (finite union)
			From measurability of $E$,
			\[
				m^*(A) = m^*(A \cap E) + m^*(A \setminus E)
			.\] 
			From measurability of $F$,
			\begin{align*}
				m^*(A \setminus E) &= m^*\left( (A \setminus E) \cap F \right)  + m^*\left( (A \setminus E) \setminus F \right) \\
				&=   m^*\left( (A \setminus E) \cap F \right)  + m^*\left(A \setminus( E \cup F) \right) 
			.\end{align*} 
			Also from the set equation $ (A \cap E) \cup \left( (A \cap F) \setminus (E \cap F) \right)  = A \cap (E \cup F)$,
			\[
				m^*(A \cap E) + m^*\left( (A \cap F) \setminus (E \cap F) \right) \ge m^*\left( A \cap (E \cup F) \right) 
			.\] 
			Thus
			\begin{align*}
				m^*(A) &= m^*(A \cap E) + m^*\left( (A \setminus E) \cap F \right) +m^*\left( A \setminus (E \cup F) \right)  \\
				&\ge m^*\left( A \cap (E \cup F) \right) +m^*\left( A\setminus (E \cup F) \right)  
			.\end{align*} 
			This is enough to show that $E \cup F$ measurable.

		\item (finite intersection) Apply deMorgan's Law to the finite union.
		\item (set exclusion) Use $E \setminus F = E \cap (\mathbb{R}\setminus F)$.
		\item (countable union)
			\begin{itemize}
				\item First assume that $E_n$ are pairwise disjoint. 
					Then
					\begin{align*}
						m^*(A) &= m^*\left( A \cap \bigcup_{1} ^p E_n \right) + m^*\left( A \setminus \bigcup_{1} ^p E_n \right) \\
						&\ge \sum_{n=1}^{p} m^*(A \cap E_n)+ m^*\left( A \setminus   \bigcup_{1} ^\infty  E_n \right)
					.\end{align*} 
					Hence
					\begin{align*}
						m^*(A) &\ge \sum_{n=1}^{\infty} m^*(A \cap E_n) + m^*\left( A \setminus \bigcup_{1} ^\infty  E_n \right) \\
						&\ge m^*\left( \bigcup_{n = 1} ^\infty (A \cap E_n) \right) + m^*\left( A \setminus \bigcup_{n = 1} ^\infty \cap E_n \right)\\
						&=  m^*\left(A \bigcup_{n = 1} ^\infty  \cap E_n \right) + m^*\left( A \setminus \bigcup_{n = 1} ^\infty \cap E_n \right)
					.\end{align*} 
					So $\bigcup_{n = 1} ^\infty E_n \in \mathcal{M}$.
					
					Now take $A = \bigcup_{n = 1} ^\infty E_n$ :
					\begin{align*}
						m^*\left( \bigcup_{n = 1} ^\infty E_n \right) &\ge \sum_{n=1}^{\infty} m^*(A \cap E_n) + 0 \\
						&=\sum_{n=1}^{\infty} m^*(E_n) + 0 
					.\end{align*}
					The other direction comes from subadditivity.
				\item In general, we let $E'_1 = E_1, E_2' = E_2\setminus E_1 \in \mathcal{M}$, etc. Now $E_n'$ are disjoint, and we reduce to the last case.
			\end{itemize}
	\end{itemize}
	
\end{proof}


\begin{definition}[sigma algebra]
	Let $X$ be a set, $A \subset 2^X$.

	$\mathcal{A}$ is an \textit{algebra} if 
	\begin{itemize}
		\item $\emptyset , X \in A$
		\item If $E, F \in \mathcal{A}$, then $E \cup F, E \cap F, E \setminus F \in \mathcal{A}$ .
	\end{itemize}

	$\mathcal{A}$ is a $\sigma$-algebra if in addition,
	\[
		E_n \in \mathcal{A}, n \in \mathbb{N} \implies \bigcup_{n = 1} ^\infty  E_n \in \mathcal{A}
	.\] 
\end{definition}

Therefore, the last theorem says that $\mathcal{M}$ is a $\sigma$-algebra.

\begin{definition}[Generated sigma algebra]
	Let $\mathcal{C}$ be a collection of subsets of $X$. The $\sigma$-algebra \textit{generated by} $X$ is the intersection of all $\sigma$-algebra containing $\mathcal{C}$. It is the smallest $\sigma$-algebra containing $\mathcal{C}$. 
\end{definition}
